%%%%%%%%%%%%%%%%%%%%%%%%%%%%%%%%%%%%%%%%%%%%%%%%%%%%%%%%%%%%%%%%%%%%%%%%%%%%%%%%%%%%%%%%%%%%%%%%
%
% CSCI 1430 Written Question Template
%
% This is a LaTeX document. LaTeX is a markup language for producing documents.
% Your task is to answer the questions by filling out this document, then to
% compile this into a PDF document.
%
% TO COMPILE:
% > pdflatex thisfile.tex

% If you do not have LaTeX, your options are:
% - VSCode extension: https://marketplace.visualstudio.com/items?itemName=James-Yu.latex-workshop
% - Online Tool: https://www.overleaf.com/ - most LaTeX packages are pre-installed here (e.g., \usepackage{}).
% - Personal laptops (all common OS): http://www.latex-project.org/get/ 
%
% If you need help with LaTeX, please come to office hours.
% Or, there is plenty of help online:
% https://en.wikibooks.org/wiki/LaTeX
%
% Good luck!
% The CSCI 1430 staff
%
%%%%%%%%%%%%%%%%%%%%%%%%%%%%%%%%%%%%%%%%%%%%%%%%%%%%%%%%%%%%%%%%%%%%%%%%%%%%%%%%%%%%%%%%%%%%%%%%
%
% How to include two graphics on the same line:
% 
% \includegraphics[width=0.49\linewidth]{yourgraphic1.png}
% \includegraphics[width=0.49\linewidth]{yourgraphic2.png}
%
% How to include equations:
%
% \begin{equation}
% y = mx+c
% \end{equation}
% 
%%%%%%%%%%%%%%%%%%%%%%%%%%%%%%%%%%%%%%%%%%%%%%%%%%%%%%%%%%%%%%%%%%%%%%%%%%%%%%%%%%%%%%%%%%%%%%

\documentclass[11pt]{article}

\usepackage[english]{babel}
\usepackage[utf8]{inputenc}
\usepackage[colorlinks = true,
            linkcolor = blue,
            urlcolor  = blue]{hyperref}
\usepackage[a4paper,margin=1.5in]{geometry}
\usepackage{graphicx}
\usepackage{fancyhdr}
\setlength{\headheight}{15pt}
\usepackage{microtype}
\usepackage{booktabs}

\frenchspacing
\setlength{\parindent}{0cm} % Default is 15pt.
\setlength{\parskip}{0.3cm plus1mm minus1mm}

\pagestyle{fancy}
\fancyhf{}
\lhead{Final Project Progress Report}
\rhead{CSCI 1430}
\rfoot{\thepage}

\date{}

\title{\vspace{-1cm}Final Project Progress Report}
\author{}

\begin{document}
\maketitle
\vspace{-1cm}
\thispagestyle{fancy}
\textbf{Important:} In your report, please\\
1) make it very clear who on the team contributed what, and\\
2) include the dates of at least \emph{two} meetings you had with your mentor.\\
\textbf{Team name: \emph{Watch Me Dance}}\\
\textbf{TA name: \emph{Kamya}}

\emph{Note:} when submitting this document to Gradescope, make sure to add all other team members to the submission. This can be done on the submission page after uploading.

\section*{Progress Report Instructions}

Before writing your progress report, you should have met with your TA and talked through your progress.

\subsection*{Team contributions}

Please describe what each of you contributed to the project so far.

\textbf{Person 1}
\begin{enumerate}
  \item Gathered a motion capture pipeline that converts RGB video into structured skeletal motion data in BVH format using a 2D-to-3D pose lifting approach.
  \item Tested 2D pose estimation frameworks including OpenPose and MediaPipe to verify compatibility with the motion capture pipeline
  \item Exploring replacing a core model with a custom-trained alternative to better align with downstream 3D pose lifting and animation requirements
\end{enumerate}

\textbf{Person 2}
\begin{enumerate}
  \item Integrated BVH motion data into Meta Animated Drawings to drive character animation to demonstrate that externally generated motion capture data can be applied to a stylized rendering system.
  \item Enabled GUI-based animation rendering on the Oscar, configuring the environment to support visualization workflows for the animation pipeline
  \item Improving motion retargeting between the BVH skeleton and the animation rig, handling differences in proportions and joint semantics, and enhancing visual coherence
\end{enumerate}

\end{document}
